\documentclass[../main.tex]{subfiles}
\begin{document}

\section{Purpose and scope of the evaluation}
\label{sec:eval-objectives}

Chapter~4 translated observed S7comm behaviour into Suricata signatures~\cite{suricata_oisf}. This chapter reports whether those signatures fire as intended when traffic from the Chapter~3 scenarios---live on the testbed or replayed from recorded captures---reaches a passive sensor. Suricata does not sit inline; it inspects mirrored copies only. The evaluation therefore answers four questions:

\begin{enumerate}
    \item \textbf{Operational impact.} Rule firings do not alter PLC state. Any consequence on the physical process would depend on operator response, which this thesis does not model.

    \item \textbf{Background alert rate.} Under legitimate engineering reads alone, how often do rules fire, and which SIDs dominate?

    \item \textbf{Attack coverage.} For each attack class in Chapter~3, does at least one SID from the associated rule group appear while that traffic is present?

    \item \textbf{Sensor capacity.} Under baseline load and under the DoS phase, does the capture path drop packets before Suricata decodes them?
\end{enumerate}

The testbed uses one OpenPLC instance, one attacker workstation, and one IDS VM (Section~\ref{sec:lab-setup}). Reported figures describe this configuration; they are not presented as universal production bounds.

\section{Evaluation methodology}
\label{sec:eval-setup}

\subsection{Measurement design}

Evaluation combines three measurements in one test campaign. Table~\ref{tab:eval-lab-setup} lists the experimental parameters.

\paragraph{Scenario detection rate (Layer~1).}
Ten attack classes were exercised: seven through live scripts on OpenPLC, block upload and download through offline replay of a reference Siemens S7comm capture~\cite{wireshark_s7comm} (the same PDU sequence documented in Section~\ref{sec:down_up_program}), and the MITM scenario through the Stuxnet simulation in Section~\ref{sec:stuxnet_mitm_sim}. A scenario \emph{meets the detection criterion} when the alert log contains at least one SID from the expected set in Chapter~4. Setup Communication (\texttt{sid:1000001}) accompanies almost every new session and is not used alone to judge attack coverage unless no more specific rule fired. The aggregate metric is
\begin{equation}
    R_{\mathrm{det}} = \frac{\#\{\text{scenarios meeting the criterion}\}}{N_{\mathrm{scenarios}}}, \quad N_{\mathrm{scenarios}} = 10.
\end{equation}
Meeting the criterion confirms that the rule set \emph{reacts} to the attack class; it does not require every optional SID in a group to fire.

\paragraph{Baseline alert density (Layer~2).}
Before attacks, the sensor observed traffic for $T = 60$\,s while the attacker performed two read-only Snap7 block-list operations. Alerts in this window are \emph{operational alerts}. Those that match legitimate protocol steps---for example Setup when Snap7 opens a session---are \emph{operational false positives} in the SOC sense: the signature matched correctly, but no attack was intended.
\begin{equation}
    \label{eq:baseline-density}
    \rho_{\mathrm{baseline}} = \frac{N_{\mathrm{alert\,baseline}}}{T} \times 3600 \quad \mathrm{(alerts/hour)}.
\end{equation}

\paragraph{Processing load (Layer~3).}
Suricata statistics were sampled at five milestones (Table~\ref{tab:eval-performance}). The primary overload indicator is \texttt{capture.kernel\_drops}. Cumulative decoded packets and alert counts show how load grows through the campaign.

\IfFileExists{../../detect/eval_results/eval_tables_en.tex}{%
  % English tables for DATN/en — data from eval_rules.py run 2026-06-25
% Regenerate Vietnamese tables with: python detect/eval_rules.py

\begin{table}[H]
    \centering
    \footnotesize
    \begin{tabular}{|l|l|}
        \hline
        \textbf{Component} & \textbf{Lab value} \\
        \hline
        IDS (Suricata) & 172.16.16.7 (lubuntu), mirror ens33 \\
        Attacker (runs scenarios) & 172.16.16.5 (ubuntu) \\
        PLC (OpenPLC :102) & 192.168.50.10 (rack 0, slot 2) \\
        HOME\_NET & \texttt{192.168.50.0/24,192.168.60.0/24} \\
        \hline
    \end{tabular}
    \caption{Rule-set evaluation lab setup}
    \label{tab:eval-lab-setup}
\end{table}

\begin{table}[H]
    \centering
    \footnotesize
    \begin{tabular}{|l|p{8.5cm}|}
        \hline
        \textbf{Metric} & \textbf{Definition in this thesis} \\
        \hline
        Scenario detection rate & Scenarios with $\geq 1$ expected SID / total scenarios \\
        \hline
        Baseline alert density & Suricata alerts during periodic read-only Snap7 only (no attack scripts) \\
        \hline
        Operational FP (baseline) & Alerts during baseline, grouped by SID (Setup, List Blocks, \ldots) \\
        \hline
        kernel\_drops & Packets dropped by the kernel on the IDS; $=0$ means no overload at lab load \\
        \hline
    \end{tabular}
    \caption{Evaluation metrics for detection and false positives}
    \label{tab:eval-metrics-def}
\end{table}

\begin{table}[H]
    \centering
    \footnotesize
    \begin{tabular}{|l|r|}
        \hline
        \textbf{Baseline} & \textbf{Value} \\
        \hline
        Observation window & 60\,s (two Snap7 reads) \\
        New alerts & 3 \\
        Alert density & 180.00 alerts/hour \\
        \hline
    \end{tabular}
    \caption{Baseline results --- legitimate operational traffic}
    \label{tab:eval-baseline}
\end{table}

\begin{table}[H]
    \centering
    \footnotesize
    \begin{tabular}{|r|r|l|}
        \hline
        \textbf{SID} & \textbf{Alert count} & \textbf{Note} \\
        \hline
        1000001 & 3 & Setup Communication --- valid Snap7 session \\
        \hline
    \end{tabular}
    \caption{SID distribution in baseline (operational false positives)}
    \label{tab:eval-baseline-sid}
\end{table}

\begin{table}[H]
    \centering
    \scriptsize
    \begin{tabular}{|l|l|c|c|l|}
        \hline
        \textbf{ID} & \textbf{Scenario} & \textbf{Alerts} & \textbf{OK?} & \textbf{SIDs observed} \\
        \hline
        \texttt{recon\_nmap} & Reconnaissance --- nmap s7-info & 6 & \checkmark & \texttt{1000001, 1000002, 1000003} \\
        \hline
        \texttt{recon\_blocks} & Reconnaissance --- list blocks (Snap7) & 1 & \checkmark & \texttt{1000001} \\
        \hline
        \texttt{write\_db} & Modify process --- Write Var DB & 5 & \checkmark & \texttt{1000001, 1000010, 1000011} \\
        \hline
        \texttt{dos\_all} & DoS --- setup / tcp / szl (all phases) & 5850 & \checkmark & \texttt{1000001, 1000002, 1000040, 1000042} \\
        \hline
        \texttt{malformed\_all} & Malformed S7comm --- all phases & 14 & \checkmark & \texttt{1000001, 1000053, 1000054} \\
        \hline
        \texttt{plc\_stop} & Start/Stop replay --- PLC Stop & 3 & \checkmark & \texttt{1000001, 1000032, 1000033} \\
        \hline
        \texttt{plc\_start} & Start/Stop replay --- PLC Start & 3 & \checkmark & \texttt{1000001, 1000030, 1000031} \\
        \hline
    \end{tabular}
    \caption{Attack detection matrix (7/7 scenarios, 100\%)}
    \label{tab:eval-detection-matrix}
\end{table}

\begin{table}[H]
    \centering
    \footnotesize
    \begin{tabular}{|l|c|c|r|r|r|}
        \hline
        \textbf{Snapshot} & \textbf{CPU (\%)} & \textbf{RAM (\%)} & \textbf{kernel\_pkts} & \textbf{drops} & \textbf{detect.alert} \\
        \hline
        baseline\_start & --- & --- & 159318 & 0 & 17467 \\
        \hline
        baseline\_end & --- & --- & 159318 & 0 & 17467 \\
        \hline
        before\_attacks & --- & --- & 159318 & 0 & 17467 \\
        \hline
        under\_dos & --- & --- & 159318 & 0 & 17467 \\
        \hline
        after\_attacks & --- & --- & 159318 & 0 & 17467 \\
        \hline
    \end{tabular}
    \caption{Suricata performance on the IDS (stats.log / top)}
    \label{tab:eval-performance}
\end{table}

\begin{table}[H]
    \centering
    \footnotesize
    \begin{tabular}{|r|r|r|l|}
        \hline
        \textbf{SID} & \textbf{Threshold} & \textbf{Window (s)} & \textbf{Description} \\
        \hline
        1000040 & 40 & 10 & TCP SYN flood :102 \\
        \hline
        1000041 & 25 & 10 & Setup Communication flood \\
        \hline
        1000042 & 25 & 10 & Read SZL flood \\
        \hline
        1000043 & 80 & 10 & S7 Job request flood \\
        \hline
    \end{tabular}
    \caption{DoS rule thresholds (ics\_dos.rules) --- must exceed HMI baseline}
    \label{tab:eval-dos-thresholds}
\end{table}
%
}{%
\begin{table}[H]
    \centering
    \caption{Evaluation tables not found --- regenerate evaluation data}
    \label{tab:eval-lab-setup}
\end{table}
}

\subsection{Reading the result tables}

Table~\ref{tab:eval-baseline} reports Layer~2: what was measured on benign traffic, the numeric outcome, and how to interpret it. Table~\ref{tab:eval-detection-summary} reports Layer~1 for each attack class: the detection objective, the number of new alerts, the expected and observed SID sets, and whether the criterion was met. Table~\ref{tab:eval-block-transfer} breaks down upload and download replay by protocol phase. Table~\ref{tab:eval-dos-thresholds} compares DoS \texttt{threshold} parameters with the maximum rates seen during baseline. Table~\ref{tab:eval-performance} reports Layer~3 at each milestone.

\section{Results and interpretation}
\label{sec:eval-results}

\subsection{Baseline behaviour}
\label{sec:eval-fp}

Table~\ref{tab:eval-baseline} summarises Layer~2. During the 60\,s window the sensor recorded three alerts, all from rule \texttt{1000001} (Setup Communication). Snap7 must establish an S7 session before listing blocks; matching Setup is therefore expected signature behaviour, not evidence of reconnaissance. No DoS rule (\texttt{1000040}--\texttt{1000043}) fired, which indicates that the configured ten-second thresholds lie above the poll rate of these benign reads (Table~\ref{tab:eval-dos-thresholds}). The scaled density $\rho_{\mathrm{baseline}} = 180$\,alerts/h is a short-window estimate; a production HMI polling continuously would require a longer baseline study before tuning alert thresholds.

\subsection{Attack scenario coverage}
\label{sec:eval-detection}

Across ten attack classes, $R_{\mathrm{det}} = 10/10 = 100\%$ (Table~\ref{tab:eval-detection-summary}). The subsections below interpret each group; the table remains the authoritative numeric record.

\paragraph{Reconnaissance.}
Network scanning produced six alerts covering Setup, generic Read SZL (\texttt{1000002}), and module-identification SZL (\texttt{1000003}). Rules for communication-status SZL (\texttt{1000004}) and block inventory (\texttt{1000005}) did not fire in the scan run, but the criterion was satisfied by the core enumeration signatures. The dedicated block-list script generated one Setup alert; List Blocks (\texttt{1000005}) did not fire on the short read, yet the criterion still holds because Setup is in the expected set for that scenario. In practice, an analyst would correlate multiple reconnaissance SIDs from the same source within a short interval rather than treat a lone Setup alert as an incident.

\paragraph{Process manipulation and MITM.}
The unauthorized Write Var to a datablock triggered both the generic Write rule (\texttt{1000010}) and the DB-area variant (\texttt{1000011}), confirming that content matches at the function byte and at the request-item structure fire on the Stuxnet-style write. The MITM scenario combined the same write path with sustained HMI polling while responses were intercepted. Rule \texttt{1000013} (elevated Ack\_Data rate from the PLC toward the HMI subnet) fired alongside the write rules, producing 47 alerts over the observation window. This pattern supports correlation---write plus abnormal read-response volume---even though \texttt{1000013} cannot prove payload tampering by itself.

\paragraph{Block upload and download.}
Upload and download were evaluated by replaying a reference capture through the IDS path so that the full Job/Ack\_Data sequence for block transfer appeared on the mirror interface. Upload generated twelve alerts: one Start Upload (\texttt{1000023}), nine Upload Block requests (\texttt{1000024}) corresponding to successive data chunks, and two End Upload (\texttt{1000025}) firings when the session closed and was confirmed. Download generated sixteen alerts: Request Download (\texttt{1000020}), thirteen firings across the bidirectional Download Block handshake (\texttt{1000021} from PLC, \texttt{1000026} from client), and two Download Ended (\texttt{1000022}). Table~\ref{tab:eval-block-transfer} maps these counts to protocol phases. The results show that opcode rules at offsets~17 and~19, together with direction filters, distinguish upload from download legs as designed in Chapter~4.

\paragraph{Denial of service.}
The DoS scenario produced 5\,850 alerts. Rate-based rules \texttt{1000040} (TCP SYN churn) and \texttt{1000042} (Read SZL flood) fired; Setup-only flood (\texttt{1000041}) and generic Job flood (\texttt{1000043}) did not, because the combined script exceeded detection through the phases that activated rather than through every threshold independently. Opcode rules such as \texttt{1000001} and \texttt{1000002} also fired on syntactically valid flooded PDUs; they provide context but do not replace the rate-based alerts. Comparing Table~\ref{tab:eval-dos-thresholds} with baseline counts explains why benign traffic stayed silent while the flood phase did not.

\paragraph{Malformed traffic.}
\label{sec:eval-malformed}
Malformed injection yielded fourteen alerts, mainly \texttt{1000053} (invalid \texttt{ROSCTR}) and \texttt{1000054} (invalid function byte). Encapsulation-level rules did not fire on every injected variant because OpenPLC often closed the TCP session before the full PDU was parsed. The criterion requires only one \texttt{100005x} alert; partial coverage across the nine fault types is expected on a lightweight stack and does not invalidate detection of the header and function-layer faults that did reach the sensor.

\paragraph{Start/Stop replay.}
Stop replay triggered \texttt{1000032} and \texttt{1000033}; Start replay triggered \texttt{1000030} and \texttt{1000031}. In both cases the generic opcode rule and the \texttt{P\_PROGRAM} string variant fired, matching the replayed payloads from Section~\ref{sec:start_stop_plc}. These rules will also match authorised maintenance unless source IP or change-window policy is applied.

\subsection{Sensor processing load}
\label{sec:eval-performance}

Table~\ref{tab:eval-performance} reports Layer~3. Decoded packets rose from approximately $1.13 \times 10^5$ at baseline end to $1.58 \times 10^5$ after all scenarios; cumulative alerts increased from 11\,583 to 17\,528. \texttt{kernel\_drops} remained zero at every snapshot, including during the DoS phase when alerts jumped from 11\,583 to 17\,446. The sensor therefore kept pace with the mirror feed at the traffic rates produced in this lab. Maximum sustainable throughput was not benchmarked; plant-scale deployment would require separate capacity testing.

\section{Discussion and limitations}
\label{sec:eval-discussion}

Several limits follow from the measurement design rather than from individual rule defects. First, detection depends on mirror visibility: traffic that bypasses the monitored segment cannot generate alerts. Second, the 60\,s baseline underestimates alert rates under continuous HMI polling; $\rho_{\mathrm{baseline}}$ should be treated as a controlled snapshot, not a production forecast. Third, intent-blind signatures on Write, Start, Stop, and block transfer cannot distinguish maintenance from attack without asset context. Fourth, the detection criterion requires only one expected SID per scenario, so partial rule-group coverage (for example three of nine malformed variants) is counted as success but limits fine-grained scoring. Fifth, block-transfer numbers come from PCAP replay rather than live OpenPLC block exchange; they validate rule logic on recorded Siemens PDUs but not path-specific quirks of the virtual PLC. Sixth, MITM rule \texttt{1000013} indicates abnormal read-response volume, not confirmed forgery. Seventh, OpenPLC's S7 stack differs from commercial CPUs, so malformed and DoS effects may not transfer one-to-one to field devices.

\label{sec:eval-dos}
\label{sec:eval-normal-impact}

\end{document}
